\documentclass{article}
\usepackage{amsmath,amssymb,mathtools}
\usepackage[margin=1in]{geometry}
\allowdisplaybreaks
\usepackage[UTF8]{ctex}
\begin{document}

\section*{具有已知污染上限的稳健动态组合优化}

\subsection*{目标与可接纳性}
所选的绩效目标是：通过一种可接受的序贯组合策略，将“典型最优典型评估”（TOTE）的遗憾程度降至最低。
\begin{equation}
\operatorname{Regret}^{\mathrm{TOTE}}(T)
:=\mathbb{E}\left[\sum_{t=1}^{T} R(S^*;\{v\})-R(S_t;\{v\})\right].
\tag{2}
\end{equation}
零售商在做出决策时所面临的约束条件包括：
\[
S_t\subseteq[N],\qquad |S_t|\le K,\qquad
S_t=\mathfrak{P}_t\bigl(\{S_\tau,i_\tau\}_{\tau\le t-1}\bigr),
\qquad t=1,2,\ldots,T,
\]
在$T$随机函数$\mathfrak{P}_1,\ldots,\mathfrak{P}_T$构成的策略中，符合以下规定的可接受策略。基准是
\[
S^*=\arg\max_{S\subseteq[N],\,|S|\le K}R(S;\{v\}).
\]
TOTE会使用典型的客户模型来评估这两种情况，即使到达的客户属于异常情况也是如此。这是该论文所选择的性能评估标准，用于进行上限分析，而不是另一种事后最佳标准。

\subsection*{输入、选择模型以及信息协议}
共有$N$个物品，以及$T$个顾客到达时间段。容量参数为$K$。各个物品的参数均满足要求。
\[
r_i\in[0,1]\quad\text{(known revenue)},\qquad
v_i\in[0,1]\quad\text{(unknown utility)},\qquad i\in[N].
\]
真正的污染参数和已知的上限有着不同的作用：
\[
\varepsilon\in[0,1),\qquad \overline{\varepsilon}\ge\varepsilon,
\]
在算法1中使用了$\overline{\varepsilon}$。所选的效用规范化方法是$v_i\in[0,1]$，而不是第7页中讨论的更大型的上界变体。

在时间$t$时，所观察到的购买行为满足以下条件：
\[
i_t\in S_t\cup\{0\},\qquad t=1,\ldots,T.
\]
购买商品$i_t\in S_t$会产生收入$r_{i_t}$；而$i_t=0$则表示没有购买任何商品，因此不会产生任何收入。对于普通客户来说，
\begin{equation}
\Pr[i_t=i\mid S_t]
=\frac{v_i}{v_0+\sum_{j\in S_t}v_j},\qquad v_0=1,
\tag{1}
\end{equation}
对于购买结果$i\in S_t\cup\{0\}$，因此对于普通客户来说，没有购买是最有可能的结果。而异常值则根据一个未知的任意分布$Q_t$进行购买，该分布在$S_t\cup\{0\}$处边缘化；$Q_t$可能会随着$t$发生变化。

源代码的假设（A1–（A2）以及定义1）明确指出了对手的情况以及零售商所掌握的信息。
\[
\phi_t\in\{0,1\},\qquad t=1,2,\ldots,T,
\]
其中，$\phi_t=1$表示异常值，而$\phi_t=0$表示典型客户。有限污染指的是……
\[
\sum_{t=1}^{T}\phi_t\le\varepsilon T. \tag{A1}
\]
自适应对手由$T$个任意可测量的函数$\mathfrak{A}_1,\ldots,\mathfrak{A}_T$组成。
\[
\begin{aligned}
\mathcal{F}_{t-1}&=\{\phi_\tau,Q_\tau,S_\tau,i_\tau\}_{\tau\le t-1},\\
\mathfrak{A}_t &: \{\phi_\tau,Q_\tau,S_\tau,i_\tau\}_{\tau\le t-1}
\longmapsto(\phi_t,Q_t),\qquad t=1,\ldots,T.
\end{aligned}
\tag{A2 / Definition 1}
\]
因此，异常顾客的到来以及选择模型可以依赖于之前的顾客、提供的商品种类，以及过去的购买行为。零售商的政策中，关于这些信息的历史记录较少。
\[
\begin{aligned}
\mathcal{G}_{t-1}&=\{S_\tau,i_\tau\}_{\tau\le t-1},\\
\mathfrak{P}_t &: \{S_\tau,i_\tau\}_{\tau\le t-1}\longmapsto S_t,
\qquad S_t\subseteq[N],\quad |S_t|\le K,
\qquad t=1,\ldots,T.
\end{aligned}
\]
如果$\phi_t=0$和$i_t$是根据(1)式实现的，且以$S_t$为条件；那么如果$\phi_t=1$，则它是根据$Q_t$来实现的。零售商不会得到异常值指标或异常值分布信息。

符号$P_t$表示客户选择的模型，而非政策函数$\mathfrak{P}_t$：
\[
P_t=\begin{cases}
\{v\},&\phi_t=0,\\
Q_t,&\phi_t=1.
\end{cases}
\]
函数$R(S;P)$表示在选择模型$P$下，提供$S$所能获得的预期收益。分别针对典型模型和估算的效用值进行计算。
\[
R(S;\{v\})=R(S;v)=\frac{\sum_{i\in S}r_iv_i}{1+\sum_{i\in S}v_i},
\qquad
R(S;\widehat v)=\frac{\sum_{i\in S}r_i\widehat v_i}{1+\sum_{i\in S}\widehat v_i}.
\]
第二个表达式是算法1中的定义性陈述；真正的效用表达式也出现在引理2的补充证明中（物理页面31）。

\subsection*{已知的上下界构造方法：算法1}
输入为 $T$、$\overline\varepsilon$、$\{r_i\}$ 和 $K$；输出为序列 $\{S_t\}_{t=1}^{T}$。初始化时，保持原有的顺序，并保留所有的数字。
\[
\widehat v^{(0)}\equiv1,\qquad
\widehat\Delta_{\overline\varepsilon}(0)=1,\qquad
\mathcal A^{(0)}=[N],\qquad
T_0=128(K+1)^2N\ln T.
\]
对于时间节点$\tau=0,1,2,\ldots$，按照指定的顺序执行以下步骤：
\[
\begin{aligned}
S_\tau^{(i)}
&=\arg\max_{S\subseteq\mathcal A^{(\tau)},\,|S|\le K,\,i\in S}
 R(S;\widehat v^{(\tau)}),
&&i\in\mathcal A^{(\tau)},\\
\gamma^{(\tau)}
&=\max_{i\in\mathcal A^{(\tau)}}R(S_\tau^{(i)};\widehat v^{(\tau)}),\\
\mathcal A^{(\tau+1)}
&=\left\{i\in\mathcal A^{(\tau)}:
R(S_\tau^{(i)};\widehat v^{(\tau)})
+2\widehat\Delta_{\overline\varepsilon}(\tau)\ge\gamma^{(\tau)}\right\},\\
n_i&=0,\qquad n_0(i)=0,
&&i\in\mathcal A^{(\tau+1)},\\
T_\tau&=2^\tau T_0.
\end{aligned}
\]
在接下来的$T_\tau$个时间周期中，随机选取$i\in\mathcal A^{(\tau+1)}$个样本，向到达的顾客提供$S_\tau^{(i)}$，观察$i_t$的情况，并更新该样本物品的计数器数值。
\[
S_t=S_\tau^{(i)},\qquad
n_i\leftarrow n_i+\mathbf1\{i_t=i\},\qquad
n_0(i)\leftarrow n_0(i)+\mathbf1\{i_t=0\}.
\]
购买的其他商品$\ell\in S_\tau^{(i)}$、$\ell\ne i$，不会计入这些计数器。在这一时期的末尾，打印出的估算值将会是……
\[
\widehat v_i^{(\tau+1)}
=\max\{1,n_i/n_0(i)\},\qquad i\in\mathcal A^{(\tau+1)}.
\]
通过以下方式定义时代污染的上限和活跃集的基数：
\[
\overline\varepsilon_\tau
=\min\{1,\overline\varepsilon T/T_\tau\},\qquad
N_\tau=|\mathcal A^{(\tau+1)}|.
\]
算法1所使用的误差上限为
\[
\widehat\Delta_{\overline\varepsilon}(\tau+1)
=\begin{cases}
1,&T_\tau<\dfrac{\overline\varepsilon T}{4(K+1)},\\[6pt]
16K(K+1)\left(
\dfrac{\overline\varepsilon_\tau}{2}
+\sqrt{\dfrac{\overline\varepsilon_\tau N_\tau\ln T}{T_\tau}}
+\dfrac{2N_\tau\ln T}{3T_\tau}
\right)
+16\sqrt{\dfrac{KN_\tau\ln T}{T_\tau}},&\text{otherwise}.
\end{cases}
\]
附带的注释将$\mathcal A^{(\tau+1)}\subseteq[N]$描述为在时代$\tau$中被探索的活跃项目；$S_\tau^{(i)}\subseteq[N]$则被描述为包含$i$的各种物品的集合；而计数器则被描述为$n_i,n_0(i)\in\mathbb N$。此外，注释还说明了其他内容。
\[
\widehat v^{(\tau)}\in[0,1]^N,\qquad
\gamma^{(\tau)}\in[0,1],\qquad
\widehat\Delta_{\overline\varepsilon}(\tau+1)\in[0,1].
\]
这些是来源的描述，以下的不一致之处被保留下来，而不是被修复。

\paragraph{此构造中存在印刷不一致的问题。}
算法1的第14行打印出了$\max\{1,n_i/n_0(i)\}$。这与所指定的范围$\widehat v^{(\tau)}\in[0,1]^N$以及引理1的补充证明（物理页面30）相矛盾。在补充证明中，当该比率不超过1时，使用$\widehat v_i=n_i/n_0$；当该比率超过1时，则使用$\widehat v_i=1$。上述更新保留了算法1中打印的最大值，而不是补充阅读中的数值。没有使用零分母的约定。同样，错误界限表达式也没有被限制在1以内，尽管其指定的范围是$[0,1]$。第9-10页的说明中提到，$\gamma^{(\tau)}$和$S_\tau^{(i)}$是通过$\mathcal A^{(\tau+1)}$计算得出的；而实际的第5-6行则使用了$\mathcal A^{(\tau)}$。伪代码给出了几何时代的长度以及一个无限制的时代循环，但没有对$T_\tau$进行四舍五入处理，也没有明确的终端部分时代规则；这里也没有添加任何相关内容。

\subsection*{链接式组合计算：第4.1节和算法2}
这是算法1所调用的计算过程，并非对顺序问题的替代目标。在调用时，候选项就是当前的$\mathcal A^{(\tau)}$，效用值就是当前的$\widehat v^{(\tau)}$，而所需的项则是$i\in\mathcal A^{(\tau)}$。第4.1节详细描述了针对特定$i\in[N]$的通用计算方法。
\begin{equation}
\max_{|S|\le K,\,i\in S}R(S;\widehat v)
=\max_{|S|\le K,\,i\in S}
\frac{\sum_{j\in S}r_j\widehat v_j}{1+\sum_{j\in S}\widehat v_j}.
\tag{4}
\end{equation}
对于任何一个$\alpha\in(0,1]$，可行性检查会判断是否存在$S\subseteq[N]$、$|S|\le K$、$i\in S$，使得…
\[
R(S;\widehat v)\ge\alpha
\quad\Longleftrightarrow\quad
\sum_{j\in S}r_j\widehat v_j
\ge\alpha+\alpha\sum_{j\in S}\widehat v_j
\quad\Longleftrightarrow\quad
\sum_{j\in S}(r_j-\alpha)\widehat v_j\ge\alpha.
\]
因为必须包含$i$，所以这个检查实际上就是在判断是否存在这样的元素。
\[
S'\subseteq[N]\setminus\{i\},\qquad |S'|\le K-1
\]
令人满意
\begin{equation}
(r_i-\alpha)\widehat v_i
+\sum_{j\in S'}(r_j-\alpha)\widehat v_j\ge\alpha.
\tag{5}
\end{equation}
附近的文本指示在$S'$中列出$(r_j-\alpha)\widehat v_j$正值最大的$(K-1)$项，并进行检查（5）。如果检查成功，则收入水平$\alpha$是可实现的；否则则不可实现。这些是计算过程中的阈值测试，而不是对每个$\alpha$提供的每种选择都施加的约束条件。

算法2提供了以下二进制搜索结果。其输入参数为$\{r_j\}_{j=1}^{n}$、$\{\widehat v_j\}_{j=1}^{n}$、必选物品$i$、$K$，以及停止精度$\delta$。其输出结果为一组$\widehat S$。
\[
|\widehat S|\le K,\qquad i\in\widehat S,
\]
最大化$R(\widehat S;\widehat v)$的值。进行初始化操作。
\[
\alpha_\ell=0,\qquad \alpha_u=1,\qquad \widehat S=\varnothing.
\]
当$\alpha_u-\alpha_\ell\ge\delta$时，设置为
\[
\alpha_{\mathrm{mid}}\leftarrow(\alpha_\ell+\alpha_u)/2.
\]
对于每个$j\ne i$，对其进行排序。
\[
\psi_j:=(r_j-\alpha_{\mathrm{mid}})\widehat v_j
\]
按降序排列，然后组合成…
\[
\Psi:=\{j\ne i:\psi_j\ge0\}.
\]
其余的打印行是：
\[
\begin{aligned}
t&:=\psi_i+\text{the $(K-1)$ $\psi_j$ in $\Psi$ with the largest values};\\
\text{if }t\ge\alpha_{\mathrm{mid}}:\quad
\widehat S&=\{i\}\cup
\{\text{the $(K-1)$ items in $\Psi$ with the largest $\psi_j$ values}\},
\qquad \alpha_\ell\leftarrow\alpha_{\mathrm{mid}};\\
\text{else}:\quad \alpha_u&\leftarrow\alpha_{\mathrm{mid}}.
\end{aligned}
\]
这里$t$是算法2中使用的标量，而不是客户到达的索引。然后循环会重复其宽度测试。

\paragraph{算法2中的打印遗漏。}
第6行将$\psi_j$仅定义为与$j\ne i$相关的内容，而第7行则使用$\psi_i$；缺失的定义并未在上方插入。第7至8行要求从$\Psi$中获取$(K-1)$项，但没有说明当可用数量不足时会发生什么情况。文中使用了正值，而伪代码中则明确使用了$\psi_j\ge0$；伪代码中的非负测试也被保留了下来。唯一明确规定的精度要求是停止测试$\alpha_u-\alpha_\ell\ge\delta$：在调用算法1时，如果没有提供$\delta$的值，或者在该有限的停止精度下无法保证精确最优性，那么就会停止计算。不过，输出结果被描述为最大化器，其初始化为空集。没有进行任何平局处理、短缺情况的处理，也没有进行任何初始化修复。

\paragraph{来源范围。}
该运筹学教材的片段来自第3节，物理页面6-8页；已知上界构造及其计算则来自第4节至4.1节，物理页面9-11页。引用的补充内容包含在提供的原始电子配套资料中。自适应未知污染策略、基于间隙的假设、数值实例以及仅用于证明的不等式，并非该选定模型的额外约束条件。

\end{document}
